\chapter{Relapsing Polychondritis}
\index{Relapsing Polychondritis}
\label{sec:Relapsing Polychondritis}

\section{Overview}
Relapsing polychondritis is a rare, immune-mediated disease characterized by recurrent, progressive inflammation and destruction of cartilaginous structures, particularly in the ears, nose, larynx, tracheobronchial tree, and joints, with an incidence of approximately 0.5--3.5 per million, equal sex distribution, and typically presenting in the fifth to sixth decade.
\section{Causes}
This condition happens because of an autoimmune attack on type II, IX, and XI collagen and other cartilage matrix components, mediated by autoantibodies, T-cell infiltration, and complement activation, with proteoglycan degradation by matrix metalloproteinases contributing to cartilage destruction.     
\section{Risk Factors}

\begin{itemize}[noitemsep]
  \item Genetic predisposition and family history
  \item Advancing age
  \item Lifestyle factors (diet, exercise, smoking, alcohol)
  \item Environmental exposures
  \item Underlying medical conditions
  \item Medication use
\end{itemize}
\section{Signs and Symptoms}

\subsection{Common symptoms}
\begin{itemize}[noitemsep]
  \item You may experience auricular chondritis (the most common presenting feature: painful, swollen, erythematous ears sparing the lobule), nasal chondritis (saddle nose deformity), laryngotracheal involvement (subglottic stenosis, hoarseness, stridor, airway collapse---the most life-threatening manifestation), costochondritis, and non-erosive seronegative arthritis, with ocular involvement (episcleritis, scleritis, uveitis) occurring in up to 60\%, audiovestibular involvement (conductive and sensorineural hearing loss, feeling of spinning) being common, and association with other autoimmune diseases (SLE, vasculitis, myelodysplastic syndrome) in 30\%. 
  \end{itemize}

\subsection{Emergency warning signs}
\begin{itemize}[noitemsep]
  \item Severe or worsening symptoms of relapsing polychondritis require immediate medical evaluation
\end{itemize}
\section{Progression and Stages}
The condition may progress through different stages of severity over time. Early stages often involve mild or subtle symptoms that may be overlooked. Moderate stages involve more noticeable symptoms affecting daily function. Advanced stages may involve significant impairment and complications.
\section{Complications}
\begin{itemize}[noitemsep]
  \item Disease progression leading to worsening symptoms
  \item Secondary infections or injuries
  \item Side effects of treatment
  \item Reduced quality of life
  \item Emotional and psychological impact
\end{itemize}
\section{Diagnosis}
Doctors usually diagnose this based on your symptoms and a physical exam, based on McAdam's criteria or Damiani's criteria, with laboratory findings being non-specific. Comprehensive medical history and physical examination Laboratory tests to assess organ function and rule out other conditions Imaging studies to visualize affected structures Specialized testing depending on the affected system Consultation with relevant specialists Monitoring of disease progression over time

\begin{itemize}[noitemsep]
  \item Comprehensive medical history and physical examination
  \item Laboratory tests to assess organ function
  \item Imaging studies to visualize affected structures
  \item Specialized testing depending on the affected system
  \item Consultation with relevant specialists
\end{itemize}
\section{Prevention}
\begin{itemize}[noitemsep]
  \item Maintain a healthy lifestyle with balanced diet and exercise
  \item Avoid known risk factors
  \item Attend regular medical check-ups for early detection
  \item Follow recommended screening guidelines
\end{itemize}
\section{Treatment Options}
Treatment options include corticosteroids (rapid symptomatic relief), dapsone (mild cases), and immunosuppressants (methotrexate, mycophenolate, cyclophosphamide) for severe disease. Medications to manage symptoms and address underlying mechanisms Lifestyle modifications and self-care measures Physical therapy and rehabilitation Surgical intervention when indicated Regular monitoring and follow-up care Patient education and support services

\begin{itemize}[noitemsep]
  \item Medications to manage symptoms and address underlying mechanisms
  \item Lifestyle modifications and self-care measures
  \item Physical therapy and rehabilitation
  \item Surgical intervention when indicated
  \item Regular monitoring and follow-up care
\end{itemize}
\section{Living with It}
\begin{itemize}[noitemsep]
  \item Follow your treatment plan as prescribed
  \item Maintain regular follow-up with your healthcare provider
  \item Make necessary lifestyle adjustments
  \item Seek support from family, friends, or support groups
  \item Stay informed about your condition
\end{itemize}
\section{Outlook}
The outlook varies from person to person, with respiratory involvement being the leading cause of death. The outlook varies depending on the specific condition, severity at diagnosis, and response to treatment. Early intervention generally leads to better outcomes. Many conditions can be effectively managed with appropriate treatment. Regular follow-up care is important for monitoring and treatment adjustment.
